\section{绪论}

\subsection{人机交互的变革和手势识别的应用}

近些年来人机交互（HCI）正经历从"有媒介"向"无缝隙"的跨越。传统交互依赖键盘、鼠标或触控屏等物理实体，限制了用户在空间上的自由度。随着普适计算的发展，交互方式正向更加自然的非接触式（Touchless）演进，手势作为人类最本能的表达语言，成为空间交互的核心。

目前主流技术包括视觉感知（如摄像头）和可穿戴感知（如IMU传感器）。视觉方案易受光照环境影响且存在严重的隐私安全风险；可穿戴方案则因其侵入性（需佩戴设备）导致用户体验受限。相比之下，多普勒雷达利用电磁波的主动探测，具备全天候工作、保护隐私及无需佩戴的显著优势。

\subsection{国内外研究进展与研究现状}

关于使用雷达对人体特征进行采集的研究早已有之\cite{lin1975noninvasive}，早期研究多基于脉冲多普勒雷达，体积大且成本高。近年来，以Google Soli\cite{lien2016soli}项目为代表，利用毫米波FMCW雷达进行微手势识别成为热点。国内清华大学、东南大学等团队在利用多普勒频率变化规律（Micro-Doppler Signature）\cite{chen2006microdoppler}提取特征方面取得了显著成果。研究重点已从简单的距离-速度探测转向了对复杂非刚体运动的精细化表征。

目前大多数研究倾向于将雷达回波转换为距离-多普勒图（RDM）或微多普勒谱图，并利用深度神经网络（CNN/RNN）进行分类。然而，这种"黑盒"模型在面对不同体型用户时鲁棒性较差。最新的研究动态（如Zhang et al.等人提出的DRS-Doppler radar sensor模型\cite{zhang2021deterministic}）开始尝试引入经典电磁散射理论，将手掌建模为米氏散射体（Mie Scatterer），试图寻找具有明确物理意义的"确定性特征"。

尽管识别准确率在实验室内已达高位，但在近场菲涅尔区（Fresnel Zone）的相位建模仍不够精确，且针对低功耗、低成本（比如5.8GHz）雷达系统的实时识别研究相对匮乏，这正是本研究的切入点。

\subsection{研究目的和研究意义}

本文旨在针对消费级雷达系统，突破传统深度学习模型对海量数据的依赖，探索一种基于物理规律引导的特征提取方法。通过建立手掌在近场环境下的米氏散射模型，精准捕捉手势动作引发的相位突变规律，实现一种低算力需求、高泛化能力、高鲁棒性的确定性手势识别方案。

本研究在理论上深化了关于非刚体在近场环境下电磁散射特性的理解，为雷达感知与手势识别的场景提供了可解释的物理支撑；并在应用上降低了手势识别对高性能运算单元的依赖，为低成本雷达芯片在智能终端的规模化部署提供了技术方案。
